\documentclass[11pt,a4paper]{article}

\usepackage[margin=2.4cm]{geometry}
\usepackage{fontspec}
\usepackage{amsmath,amssymb}
\usepackage{booktabs,longtable,array}
\usepackage{xcolor}
\usepackage{hyperref}

\setmainfont{STSong}
\setsansfont{Hiragino Sans GB}
\setmonofont{Menlo}
\XeTeXlinebreaklocale "zh"
\XeTeXlinebreakskip = 0pt plus 1pt

\hypersetup{
  colorlinks=true,
  linkcolor=blue!55!black,
  citecolor=blue!55!black,
  urlcolor=blue!55!black,
  pdftitle={PiBE-FAST BEFA 技术报告}
}
\setlength{\parindent}{2em}
\setlength{\parskip}{0.35em}
\renewcommand{\arraystretch}{1.18}
\renewcommand{\contentsname}{目录}
\renewcommand{\abstractname}{摘要}
\renewcommand{\refname}{参考文献}
\renewcommand{\appendixname}{附录}
\renewcommand{\tablename}{表}

\title{\textbf{PiBE-FAST：平衡、眼睛、面部、手臂、言语模块及研究融合层的技术设计}}
\author{项目技术报告}
\date{实现版本快照：2026 年 8 月}

\begin{document}
\maketitle

\begin{abstract}
PiBE-FAST 是一个面向树莓派边缘部署的急性脑卒中警示体征多模态辅助筛查研究原型。本报告说明系统中平衡（B）、眼睛（E）、面部（F）、手臂（A）、言语（S）五个模块及质量门控特征融合层的设计与实现。系统融合 MoveNet 人体关键点、MediaPipe Face Landmarker 面部及虹膜关键点、本地麦克风音频、离线 whisper.cpp 转写、标准化短时任务、多帧时序统计、质量门控与保守规则判定。言语模块包含两层：相对个人声学基线的自然语音变化监测，以及固定句主动确认。主动层还会记录一个基于官方普通话构音障碍语音库（MDSC）训练的 90 维 log-Mel 影子表征，但不改变 S 或紧急程度决策。融合层把已完成的 B/E/F/A/S 测量转换为固定 28 维研究向量，同时保留质量与缺失信息；它不产生诊断，也不替代已有分项或紧急程度决策。设计上优先采用本人或照护者报告的突发症状，明确区分有效阴性与数据不足，同时不把工程阈值解释为脑卒中概率。当前软件实现了可复现的测量流程和可测试的特征定义，但尚未证明临床诊断效度。出现突发神经系统症状时，不得因等待本系统检测而延误急诊评估。
\end{abstract}

\textbf{关键词：}BE-FAST；脑卒中警示体征；树莓派；姿态估计；面部关键点；眼动；言语分析；边缘计算；多模态融合。

\tableofcontents
\newpage

\section{项目背景与研究目标}

\subsection{临床与应用背景}

BE-FAST 在传统 FAST 的基础上加入平衡和视觉警示体征，目的是减少传统口诀可能遗漏的脑卒中表现\cite{aroor2017befast}。PiBE-FAST 将其中一部分转化为适合本地边缘设备运行的标准化摄像头辅助流程，面向家庭、社区、教学展示和技术研究等场景，帮助使用者完成时间较短、步骤明确的筛查。

本项目不能诊断或排除脑卒中。普通摄像头无法观测许多重要临床体征，摄像头任务结果正常也不能证明没有急性脑卒中。如果突然出现面部歪斜、单侧无力、言语不清、视觉异常或失去平衡，应立即联系急救服务，不应等待系统完成检测。

\subsection{技术目标}

BEFA 实现具有以下目标：

\begin{enumerate}
  \item 使用普通 RGB 摄像头获取有解释性的测量结果，并在本地完成推理；
  \item 通过标准化任务提高时序比较和左右侧比较的可解释性；
  \item 对低质量或未完成的测量输出“数据不足”，而不是误判为正常；
  \item 为每项结果保留特征值、质量指标和判定原因；
  \item 通过按阶段调度模型，使系统适合树莓派计算资源。
\end{enumerate}

\subsection{报告范围}

本报告详细讨论 B、E、F、A、S 及其特征级融合。发病时间（T）只在影响紧急程度的部分说明。下文中的参数来自实现快照时的 \texttt{app/befast/config.py}、\texttt{app/speech\_audio.py} 和 \texttt{app/passive\_speech.py}。除非另有说明，所有数值均为需要在目标硬件与目标人群中重新验证的工程参数。

\section{需求分析与系统架构}

\subsection{功能需求}

系统允许 B、E、F、A 独立选择和重复检测。每个项目提供准备说明、人物位置与动作质量检查，并在正式采集前倒计时。同一项目重复检测时保留每一次报告，不覆盖历史结果。自动项目可以跳过，但“已跳过”不能解释为阴性。B 和 E 首先采集人工症状信息，因为本人或照护者报告的突发失衡或视觉异常不能被摄像头结果覆盖。

\subsection{非功能需求}

系统应支持本地运行、较低计算负载、安全降级和数据最小化。模型缺失或画面不可用时仍应提供原始预览，同时将相应项目标记为数据不足。浏览器主要承担显示和控制，推理运行在 Python 服务主机。视频输入可以来自主机连接的 CSI/USB 摄像头，也可以来自当前浏览器设备上传的前置摄像头帧。

\subsection{处理架构}

\begin{center}
\begin{tabular}{c}
\toprule
主机 CSI/USB 摄像头、浏览器前置摄像头和本地麦克风 \\
$\downarrow$ \\
图像采集、缩放、方向处理和输入源路由 \\
$\downarrow$ \\
\begin{tabular}{ccc}
MoveNet & MediaPipe Face & 本地音频/whisper.cpp \\
B、A & E、F & S
\end{tabular} \\
$\downarrow$ \\
时序特征、动作完成门控与质量门控 \\
$\downarrow$ \\
BEFA 规则引擎与结构化结果 \\
$\downarrow$ \\
Flask API、React 引导界面、历史报告与研究日志 \\
\bottomrule
\end{tabular}
\end{center}

后端技术包括 Python、Flask、OpenCV 或 Picamera2、NumPy、MediaPipe、用于 MoveNet 的 TensorFlow Lite 解释器、树莓派 ALSA \texttt{arecord} 和本地 whisper.cpp。前端使用 React 和 Vite。个人基线和研究数据保存在本地结构化文件中。

\subsection{阶段化模型调度}

同时持续运行姿态和面部模型会增加计算负载和设备温度，因此 PiBE-FAST 只启用当前阶段所需的模型。待机阶段默认约以 2 FPS 运行 MoveNet；E/F 阶段暂停 MoveNet，并约以 5 FPS 运行 Face Landmarker；A 和主动 B 阶段运行 MoveNet，并暂停人脸模型；固定句 S 录音与分析期间暂停两套视觉模型。主动 S 占用麦克风前会先暂停被动自然语音监测，避免两个采集进程竞争，结束后再恢复。预览画面与实际模型推理频率相互独立。

\section{公共测量与安全方法}

\subsection{关键点与置信度过滤}

B 和 A 使用 MoveNet 人体关键点，默认最低关键点置信度为 0.35。E 和 F 使用 MediaPipe Face Landmarker 提供的 478 点面部网格、虹膜中心、表情 blendshape，以及用于头姿检查的 canonical-face transformation。只有当前特征所需关键点均存在并通过几何质量检查时，该帧才会作为有效样本。

\subsection{尺度与方向归一化}

人体距离使用肩宽归一化，以降低相机距离变化的影响。面部几何以双眼外眼角连线进行滚转校正，并使用双眼外眼角距离归一化。眼睛位置表示为虹膜在当前眼裂中的相对位置。这些处理提高了不同采集条件下的可比性，但无法完全消除透视、镜头畸变和头部深度旋转。

\subsection{多帧时序统计}

系统避免使用单帧作医学解释。实现根据动作阶段分别保存样本，并使用中位数、绝对中位差（MAD）、四分位距（IQR）、均方根变化、百分位范围和重复试次。这可以降低孤立关键点误差的影响，并证明要求的动作在一段时间内持续成立。

\subsection{结果状态语义}

\begin{table}[htbp]
\centering
\caption{各模块的公共结果状态。}
\begin{tabular}{p{0.20\linewidth}p{0.70\linewidth}}
\toprule
状态 & 含义 \\
\midrule
\texttt{checking} & 正在采集或处理。 \\
\texttt{positive} & 人工报告症状或预定义可见表型满足当前工程规则。 \\
\texttt{negative} & 动作和质量门控均通过，且未检出目标可见表型。 \\
\texttt{insufficient} & 可见性、稳定性、动作完成度、个人基线或模型能力不足。 \\
\texttt{skipped} & 使用者主动跳过，不作阴性推断。 \\
\bottomrule
\end{tabular}
\end{table}

区分 \texttt{negative} 与 \texttt{insufficient} 是系统的核心安全设计。手腕缺失、虹膜未跟踪或没有完成微笑，都不能静默地变成“正常证据”。

\section{B——平衡模块}

\subsection{目标与使用流程}

B 首先询问本人或照护者是否观察到突发失衡、协调困难或步态异常。报告新发异常时直接形成警示，不要求使用者冒险站立。如果没有报告异常，并且确认附近有支撑、站立安全，系统才使用摄像头测量当前安静站立状态，并与个人基线比较。

自动输出被明确定义为“摄像头躯干运动学代理”，不是力台压力中心（COP）、真实全身质心、临床 lateropulsion 量表或疾病诊断。Lateropulsion 涉及相对重力的身体方向以及主动推挤、抗拒纠正等临床行为，正面单目视频不能完整观测\cite{dai2021lateropulsion,dai2022lateropulsion}。

\subsection{输入关键点与单帧特征}

MoveNet 输入点包括左右肩、左右髋和左右踝。设 $(x_s,y_s)$ 为肩中心，$(x_h,y_h)$ 为髋中心，$x_a$ 为双踝中点，$x_t$ 为肩中心与髋中心的中点，$w_s$ 为肩宽，则：

\begin{align}
\theta_{\mathrm{trunk}} & = \operatorname{atan2}(x_s-x_h,\,y_h-y_s), \\
p_{\mathrm{support}} & = \frac{x_t-x_a}{w_s}.
\end{align}

前者转换为角度，后者是躯干相对双踝中点的肩宽归一化横向位移。后者有意不命名为 COP。卒中研究观察到负重不对称、姿势不稳与摆动之间的关联，但这不能证明二维踝部参照等价于力台测量\cite{kamphuis2013asymmetry,yeung2014kinect}。

\subsection{窗口级特征}

经过 1.5 秒准备期后，系统采集标称 30 秒安静站立窗口，并从有效帧序列计算：

\begin{itemize}
  \item 躯干侧倾角中位数；
  \item 归一化躯干支撑偏移中位数；
  \item 横向摆动 RMS；
  \item 横向摆动第 5--95 百分位范围；
  \item 横向路径长度和平均速度；
  \item 躯干角度围绕中位数的 RMS。
\end{itemize}

30 秒重复试次与静态平衡可靠性研究中的常见设计一致，但本项目的低频摄像头代理不能与力台指标互换\cite{aryan2023forceplate,ruhe2010cop}。

\subsection{个人基线与稳健比较}

系统默认使用 5 个合格窗口建立个人基线。对基线特征 $x$，实现计算：

\begin{align}
c_x &= \operatorname{median}(x_1,\ldots,x_n), \\
s_x &= \max\left(1.4826\,\operatorname{MAD}(x),\frac{\operatorname{IQR}(x)}{1.349}\right), \\
z_x &= \frac{|x_{\mathrm{current}}-c_x|}{s_x}.
\end{align}

躯干方向采用双侧变化，摆动平均速度只把增加作为异常方向。分数达到 3.5 表示统计过程变化\cite{iglewicz1993outliers}，并不是卒中临床 cutoff。如果 MAD 和 IQR 均为零，而之后数值发生变化，则分数没有定义，系统返回数据不足，不使用人为噪声下限强制产生阳性或阴性。

\subsection{质量门控与判定}

每个窗口至少需要 30 个有效样本，有效帧比例不低于 0.75。在 5 个基线窗口收集完成前，结果为 \texttt{personal\_balance\_baseline\_calibrating}。基线可用后，结果可能为持续躯干方向变化、横向摆动速度增加，或未见相对个人基线的持续变化。可见倾斜方向只报告为画面方向，不转换为脑损伤侧别。

\subsection{局限性}

单目姿态估计不能测量前庭症状、共济失调、前后方向摆动、主动推挤、抗拒纠正或真实负重分布。相机位置、地面、助行器、骨科疾病和既往残疾都会改变特征。与个人基线一致不能排除脑卒中。

\section{E——眼睛模块}

\subsection{目标与症状优先原则}

E 首先记录突发视力下降、黑蒙、视野缺损、复视或持续凝视偏向。本人报告新发视觉异常时直接产生 E 警示，摄像头结果不能覆盖。未报告异常时，摄像头任务辅助检查大幅静息共轭偏向、可重复的双眼方向性终点减弱和明显不共轭响应。它不检查视力、视野、眼底，也不能覆盖所有眼球运动障碍。

急性脑卒中的共轭凝视偏向和定量眼动异常已有相关研究\cite{singer2006gaze,kudo2021eyemovement}，但影像测角或研究级眼动仪得到的数值不能直接迁移到低帧率普通摄像头。

\subsection{关键点与标准化目标序列}

系统使用右虹膜中心 468、左虹膜中心 473、右眼角 33/133、左眼角 362/263，以及外眼角 33/263 完成尺度与滚转校正。MediaPipe Iris 提供了相关虹膜关键点表示\cite{bazarevsky2020mediapipeiris}。

使用者首先自然直视。随后页面交替显示中心和侧方目标，使每个侧方终点都有一个紧邻的中心终点作为参照。左右目标各重复三次，每个目标持续 2.0 秒，切换后前 0.5 秒不进入统计。页面根据屏幕物理宽度和观看距离，将侧方目标放置在约 $\pm15^{\circ}$，并记录经视口限幅后的实际角度。

\subsection{眼位与响应特征}

单眼虹膜在眼裂内的水平位置为：

\begin{equation}
g_x=\frac{x_{\mathrm{iris}}-x_{\mathrm{corner,min}}}{w_{\mathrm{eye}}}.
\end{equation}

配对侧方响应是侧方终点与前一个中心终点的有符号差，其稳健信噪比为：

\begin{equation}
\mathrm{SNR}=\frac{r}{1.4826\left(\mathrm{MAD}_{\mathrm{centre}}+\mathrm{MAD}_{\mathrm{lateral}}\right)}.
\end{equation}

系统根据左右终点总体排列自动识别摄像头镜像方向。每只眼、每个方向使用三次响应的中位数，并允许一个离群试次；至少两次需要方向正确并通过 SNR 门槛。

响应门控通过后，系统计算双眼方向不对称、每只眼的方向份额、双眼共轭误差，以及结合本轮实际目标角度估计的静息眼位偏差。系统有意不输出扫视潜伏期、峰值速度、平滑追踪增益或眼震，因为约 5 FPS 不适合这些动态指标，采样频率会显著影响扫视分析\cite{bahill1975mainsequence}。

\subsection{质量门控}

每个试次至少需要 5 个有效终点样本，有效比例不低于 0.60。每只眼实际宽度至少为 24 像素；终点 MAD 不高于 0.08；至少 80\% 有效帧必须包含头姿；单轮 pitch、yaw 和 roll 变化均不超过 $8^{\circ}$；最接近中位数的重复试次相对误差不超过 1.00。任一条件失败都返回 \texttt{insufficient}。

\subsection{判定顺序}

判定按以下顺序执行：

\begin{enumerate}
  \item 四个“眼睛--方向”响应全部失败时，说明没有证明完成目标跟随，结果为数据不足；
  \item 双眼在同一方向持续失败，提示双眼方向性凝视受限；
  \item 单眼特定方向持续失败，提示可能不共轭凝视受限；
  \item 双眼同向静息偏差达到 $12^{\circ}$，提示静息共轭凝视偏向；
  \item 双眼平均方向不对称达到 0.45，提示方向性眼动减弱；
  \item 最大归一化共轭误差达到 0.35，提示可能核间性不共轭；
  \item 其余情况报告未见可重复的可见眼动终点异常。
\end{enumerate}

响应 SNR 门槛为 3.5。$12^{\circ}$、0.45 和 0.35 都是需要按设备和人群校准的研究候选值，不能表述为经过验证的临床界限。

\subsection{局限性}

眼镜、反光、眼睑遮挡、眼型、面部尺寸、头部代偿和相机角度会影响跟踪。该任务不能识别视野缺损、视网膜病变、短暂黑蒙或许多复视病因。阴性终点结果只适用于本轮质量合格条件下所检查的可见表型。

\section{F——面部模块}

\subsection{目标与检测流程}

F 检测从中性表情到微笑时的下半脸运动不对称。与单张静态图像相比，动态个人基线可以降低天然脸型不对称和固定相机滚转的影响，但不能完全消除。使用者先保持约 2 秒中性表情，再自然微笑约 3 秒。

\subsection{特征定义}

关键点 61 和 291 分别表示右、左嘴角，双眼外眼角 33 和 263 定义滚转校正后的面部坐标和眼距尺度。MediaPipe 的 \texttt{mouthSmileLeft} 与 \texttt{mouthSmileRight} 提供第二类表情表示。

各阶段使用中位数。设 $d_n$ 和 $d_s$ 分别为中性与微笑阶段的归一化嘴角垂直差，则基线校正后的几何特征为：

\begin{equation}
\Delta d=d_s-d_n.
\end{equation}

表情特征为微笑阶段左右 smile blendshape 中位数之差。系统保留几何特征和表情模型特征的可解释性，不将其融合为不可解释的卒中概率。

\subsection{质量门控与判定}

中性和微笑阶段各至少需要 5 个有效样本，总有效帧比例不低于 0.50。左右最大微笑激活低于 0.22 时，说明没有证明完成足够微笑，结果为数据不足。当 $|\Delta d|\geq0.075$ 个眼距单位，或左右 blendshape 激活差绝对值不低于 0.24 时，结果为阳性。占主导的有符号特征用于报告可能受影响侧；其余情况报告未见明确微笑不对称。

\subsection{局限性}

既往面瘫、天然面部不对称、牙科或颌面疾病、胡须、口罩、光照和大幅转头均会影响测量。侧别标签描述可见面部特征，不用于定位脑部病灶。当前阈值需要外部验证。

\section{A——手臂模块}

\subsection{目标与标准化动作}

A 首先验证使用者是否完整完成双臂“抬起--保持--放下”，然后评估保持阶段的持续高度差和单侧下落。使用者先观看完整动作，在正面坐姿下保证双肩、双肘、双腕和双髋可见。至少预览 2 秒后，双臂自然下垂稳定约 1.2 秒即可自动进入实心 3--2--1 倒计时。正式采集依次要求：

\begin{enumerate}
  \item 起始时双臂自然下垂；
  \item 在 8 秒内将双臂横向抬至肩高；
  \item 双臂有效保持 5 秒；
  \item 在 5 秒内将双臂放下。
\end{enumerate}

每个要求的姿势需要持续约 0.45 秒。任一阶段失败都会清除半轮数据并要求重试。只有保持阶段样本用于不对称判定。

\subsection{动作完成特征}

系统将 MoveNet 的肩、肘、腕和髋关键点转换为肩宽归一化特征。准备姿势和抬臂完成门控使用手腕相对肩部高度、肘关节角度和横向伸展。默认抬臂要求包括：手腕高度容差 0.42、肘角不低于 $125^{\circ}$、横向伸展不低于 0.45。

\subsection{不对称时序特征}

设 $y_L(t)$、$y_R(t)$ 为以向下为正的左右手腕相对肩部归一化高度。系统计算保持阶段双臂高度差中位数，并将样本分为三段，比较第一段和最后一段：

\begin{align}
\delta_L &= \operatorname{median}(y_L^{\mathrm{final}})-\operatorname{median}(y_L^{\mathrm{initial}}),\\
\delta_R &= \operatorname{median}(y_R^{\mathrm{final}})-\operatorname{median}(y_R^{\mathrm{initial}}),\\
\Delta_{\mathrm{drift}} &= \delta_L-\delta_R.
\end{align}

按照内部坐标定义，正值表示左侧出现更大的向下位移。

\subsection{质量门控与判定}

只有完整动作通过后才进入分类。保持阶段至少需要 20 个有效样本；实现同时记录有效帧比例，配置中的名义质量目标为 0.55。双臂归一化持续高度差绝对值达到 0.30 时提示高度不对称；$|\Delta_{\mathrm{drift}}|\geq0.22$ 时提示不对称下落。两者同时阳性时，使用相对各自阈值幅度更大的特征确定报告侧别。均未达到时，结果为 \texttt{no\_clear\_arm\_asymmetry}。

代码库还包含 Toronto 代偿模型和 IntelliRehab 动作质量模型。它们的输出统一使用 \path{shadow_} 前缀，仅作为研究观察，不改变用户可见的 A 判定，因为其原始任务、传感器和本项目急性双臂筛查协议并不一致。

\subsection{局限性}

疼痛、肩部疾病、既往无力、关节活动受限、衣物、轮椅坐姿和遮挡会影响动作完成。单目二维姿态无法完整描述前后深度、前臂旋前、肌张力或精细远端无力。当前动作协议和阈值需要与目标摄像头下的临床专家标注比较验证。

\section{S——言语模块}

\subsection{双层设计与医学角色}

言语模块包含两个有意分离的层次。被动层从自然语音窗口中检测相对同一人声学基线的持续变化，不运行语音识别，只能建议进入固定句确认。主动层录制页面显示的固定句，在本地离线转写，并评估录音质量、复述差异、语速与停顿。两层都不能诊断构音障碍或脑卒中。无论自动结果如何，突发言语困难都属于急症警示。

卒中后构音障碍可能影响言语时序、发声及构音/共振，相关临床和声学研究为特征域提供了依据\cite{decock2021acute,kim2011acoustic,mou2018mandarin,ge2021vowels,wong2022tdcs,sanguedolce2025stroke,jyothi2026stroke}。但这些研究使用临床评估、受控元音、朗读、图片描述或带标注的连接语音，并没有直接验证家庭长期监听、本项目固定句或下述阈值。

\subsection{被动自然语音监测}

被动监测默认采集 8 秒、16 kHz 单声道窗口，每次完成采集后间隔 1 秒。每个窗口首先通过音频质量门控：时长至少 2.0 秒、RMS 不低于 $-42$ dBFS、削波比例不高于 0.02、有效发声时长至少 1.5 秒。无效或静音窗口不会进入个人基线。

轻量本地特征分为三个证据域：

\begin{longtable}{p{0.22\linewidth}p{0.68\linewidth}}
\toprule
证据域 & 特征 \\
\midrule
时序 & 停顿比例、平均停顿时长、基于能量峰的音节核速率代理 \\
发声 & 强度范围、F0 中位数、F0 半音标准差、局部 jitter、局部 shimmer、谐噪比 \\
构音/共振 & 从自由语音 LPC 帧估计的 F1、F2 四分位距 \\
\bottomrule
\end{longtable}

这些特征是适合边缘计算的近似量。音节核速率不等于强制对齐构音速率；基频和音质估计不等于 Praat 或 eGeMAPS；自由语音共振峰 IQR 也不等于元音空间面积、VAI 或 FCR。

系统默认使用 6 个有效窗口建立个人基线，最多保存 48 个基线历史窗口；一项特征至少需要 3 个基线观测才能比较。对特征 $f$，变化分数为：

\begin{equation}
z_f=\frac{|f_{\mathrm{current}}-\operatorname{median}(f_{\mathrm{baseline}})|}
{\max\left(1.4826\operatorname{MAD}(f_{\mathrm{baseline}}),s_{f,\min}\right)},
\end{equation}

其中 $s_{f,\min}$ 是按特征定义的工程最小尺度。一个证据域中至少一项可用特征满足 $z_f\geq3.5$，该域才记为变化；同一窗口至少两个域变化，才投一张异常票；最近 5 个有效窗口中至少 3 票异常，页面才建议进行主动固定句 S。单一特征、单一域或单个窗口都不能触发建议。

API 将被动输出标记为 \path{change_detection_trigger_only}，同时明确返回“尚未临床验证”，并提供证据版本、变化域、分域分数和投票历史。监测器不验证说话人身份，因此电视、其他说话人、呼吸道疾病、麦克风移动和正常个体内变化都可能影响结果。

\subsection{主动固定句采集}

使用者启动 S 后，被动采集先暂停，服务端录制约 7 秒、16-bit、16 kHz 单声道 PCM。树莓派使用 ALSA \texttt{arecord}，macOS 开发环境使用 FFmpeg/AVFoundation。默认提示句为：

\begin{quote}
中文：今天天气很好，我们一起去公园散步\\
英文：The sky is blue and we are walking in the park
\end{quote}

异步服务依次进入录音、分析和结果就绪状态。本地 whisper.cpp 按选定语言转写 WAV。麦克风工具缺失、模型或可执行程序缺失、识别失败或转写为空时，结果均为 \texttt{insufficient}，绝不会把识别能力缺失当作正常。

\subsection{音频质量与主动特征}

WAV 被划分为 20 ms 帧。系统从完整波形计算 RMS 和削波比例；语音活动阈值综合使用帧 RMS 第 20 百分位、$-42$ dBFS 绝对下限，以及相对整段 RMS 的上限。有效发声时长来自活动帧数量；停顿比例只在第一个至最后一个活动帧之间计算。

质量门控要求录音时长至少 2.0 秒、有效发声至少 1.0 秒、RMS 不低于 $-42$ dBFS、削波比例不高于 0.02。转写归一化后，字符错误率定义为：

\begin{equation}
\mathrm{CER}=\frac{D_{\mathrm{Lev}}(t_{\mathrm{expected}},t_{\mathrm{recognized}})}
{|t_{\mathrm{expected}}|},
\end{equation}

语速为归一化识别字符数除以有效发声秒数。

\subsection{主动判定顺序}

当前规则按以下顺序执行：

\begin{enumerate}
  \item 录音过短、过小、削波严重或没有足够发声时返回 \texttt{insufficient}；
  \item 识别器不可用或转写为空时返回 \texttt{insufficient}；
  \item $\mathrm{CER}\geq0.35$ 时返回 \path{speech_content_mismatch}；
  \item 停顿比例 $\geq0.55$、语速低于每秒 1.0 字符或高于每秒 8.0 字符时返回 \path{speech_delivery_irregular}；
  \item 其余情况返回 \path{no_clear_speech_abnormality}。
\end{enumerate}

这些规则检查固定句是否被可识别、连续地复述。ASR 差异不等于构音障碍：方言、口音、听力、语言熟练度、背景噪声、认知、失语和识别器本身的错误都可能造成差异。反过来，使用者也可能成功复述固定句，但仍存在临床上重要的言语或语言症状。

\subsection{MDSC 构音障碍影子表征}

主动层另外加载 \texttt{models/mdsc\_dysarthria\_v1.json}。MDSC 包含来自 21 名构音障碍和 25 名对照普通话说话人的 18,630 条 16-kHz 录音\cite{gao2024mdsc}。本地审计将每条转写与 WAV 一一匹配并保留全部录音，其中 18,616 条为单声道、14 条为立体声；后者由训练与运行共享的读取器做通道平均。该数据表示慢性构音障碍表型和受控手机麦克风采集，不是急性脑卒中或个体内急性变化。

表征使用 25 ms 帧长、10 ms 帧移的 40-bin log-Mel 能量。每个频带的均值和标准差贡献 80 维；时长、RMS、有声与停顿比例、平均停顿时长、能量变化、谱流、过零率和谱质心均值/标准差再贡献 10 维。标准化逻辑回归因此使用 90 维特征，并导出为运行时只需 NumPy 的轻量 JSON。

使用固定随机种子、在每个类别内按说话人分配 33 名训练、7 名验证和 6 名完全持出测试参与者，任一说话人都不会跨分区。40 名开发说话人执行五折说话人分组预测；部署系数只用 33 名训练说话人拟合，决策阈值 \texttt{0.807859} 只由 7 名验证说话人的平均概率选择。测试说话人只在最终报告中使用一次。表~\ref{tab:mdsc-results-zh} 同时给出录音级和说话人级结果。

\begin{table}[htbp]
\centering
\small
\caption{MDSC v1 内部评估。测试说话人只有 3 名阳性和 3 名阴性。}
\label{tab:mdsc-results-zh}
\begin{tabular}{lccccc}
\toprule
评估单位 & ROC-AUC & 敏感度 & 特异度 & 平衡准确率 & Brier \\
\midrule
开发集折外录音（$n=16{,}200$） & 0.8241 & 0.7956 & 0.6655 & 0.7306 & 0.1942 \\
开发集折外说话人（$n=40$） & 0.9015 & 0.9444 & 0.5909 & 0.7677 & 0.1272 \\
完全持出录音（$n=2{,}430$） & 0.9573 & 0.6444 & 0.9926 & 0.8185 & 0.0983 \\
完全持出说话人（$n=6$） & 1.0000 & 0.6667 & 1.0000 & 0.8333 & 0.0577 \\
\bottomrule
\end{tabular}
\end{table}

说话人级 Bootstrap 区间以说话人而不是录音为重采样单位。测试说话人只有 6 名，持出敏感度区间为 0.0--1.0，说话人级 ROC-AUC 1.0 同样不稳定。这些数字只是内部研究基线，不是临床性能。开发机上 100 次真实 WAV 推理的中位时间为 0.96 ms、第 95 百分位为 1.33 ms；这不是树莓派 M5 性能。

运行时模型向主动 S 报告增加 \path{shadow_dysarthria_probability}、冻结阈值、模型版本和明确的 \path{dysarthria_representation_not_acute_stroke} 角色。模型缺失、校验失败或推理失败时，Whisper、质量门控、原有 S 结果、融合和紧急逻辑保持不变。该表征不参与被动自然语音投票。目标麦克风 M5 仍未完成：需在安静 20 cm、安静 1 m 和家庭噪声条件下前瞻性测试冻结模型，执行 30 分钟稳定性运行，报告延迟与内存，且不能根据影子参与者反复调整阈值。

\subsection{发病时间、存储与融合接口}

主动请求会记录言语问题是否新发或突然发生，并可保存首次发现时间。S 阳性且新发/突然发生时，现有 T 逻辑将单项报告判为紧急。S 报告来源标记为 \path{microphone+offline_asr}。

日常被动 WAV 在特征提取后删除，被动基线只保存数值特征，不保存日常原始语音。主动临时音频在报告处理后删除；如果配置的事件历史需要保留阳性事件，则必须遵循项目明确的数据保留和知情同意策略。28 维融合层只使用已完成主动 S 报告中的三个指标：CER、每秒字符数和停顿比例；被动触发投票不会写入模型向量。

\subsection{局限性与验证要求}

被动监测没有说话人验证，也没有在家庭开放麦克风条件下完成临床验证。主动层部分依赖 ASR，评估的是固定句而不是言语语言治疗师检查。MDSC 验证不能弥合这一差距，因为其阳性标签是慢性构音障碍而不是急性脑卒中，且采集协议与目标麦克风不同。验证应纳入急性脑卒中、卒中类似疾病、既往构音障碍或失语、健康对照，并覆盖方言、年龄、性别、麦克风距离、多人说话、电视和呼吸道疾病。数据必须按受试者划分，参考标准应包含盲法言语语言评估、神经科诊断和发病时间；同时报告灵敏度、特异度、PPV、NPV、校准、置信区间、数据不足率和基线漂移。

\section{质量门控特征融合}

\subsection{目的与安全边界}

新增的 \texttt{befast-feature-fusion-v1} 层从已经完成的分项报告中构建固定长度、可检查的早期融合载荷，用于日志记录、探索分析和未来按受试者训练模型。它明确声明为 \path{research_only_no_decision} 模式，不改变 B/E/F/A/S 的分项结果，不输出脑卒中概率，也不替代现有发病时间安全规则或 \path{screening_decision}。

这一分离是有意的：一个便于计算的聚合值，只有在预测目标、数据集、校准、工作点和外部效度都明确后，才可能具有临床意义。仅由人工报告产生的阳性结果仍会在普通报告中发挥安全作用，但在传感器融合向量中标记为缺失，不会被错误表示成“传感器测量值为零”。

\subsection{可用性、原始特征与固定维数}

一个分项只有同时满足以下条件才视为可用于融合：报告状态为 \texttt{positive} 或 \texttt{negative}、指标映射非空、结果质量大于零。\texttt{insufficient}、\texttt{skipped}、等待状态以及只有人工报告而没有传感器指标的项目，都属于传感器测量不可用。

原始向量按固定顺序包含 18 个连续特征：

\begin{longtable}{p{0.10\linewidth}p{0.80\linewidth}}
\toprule
分域 & 连续特征 \\
\midrule
B & 躯干方向变化分数；横向摆动变化分数 \\
E & 左右眼活动范围；两眼范围不对称；双眼方向不对称；最大共轭误差；静息共轭偏向角度 \\
F & 嘴角变化；左右微笑分数差；总体微笑强度 \\
A & 双臂高度差；左臂下落；右臂下落；左右下落差 \\
S & 字符错误率；每秒字符数；停顿比例 \\
\bottomrule
\end{longtable}

缺失或非有限特征值会以零占位，但仅在明确保存相应质量和缺失字段之后执行。最终模型向量为：

\begin{equation}
\mathbf{x}_{\mathrm{model}}=
\left[
\mathbf{x}_{\mathrm{gated}}^{(18)},
q_B,q_E,q_F,q_A,q_S,
m_B,m_E,m_F,m_A,m_S
\right],
\end{equation}

其中 $q_c\in[0,1]$ 为分项质量，$m_c\in\{0,1\}$ 为缺失掩码，因此总维数为 $18+5+5=28$。这样可以区分“数值确实为零”和“该项没有可用测量”。

\subsection{质量门控}

每个特征属于一个分域 $c(j)$，门控后的数值为：

\begin{equation}
x_{j,\mathrm{gated}}=q_{c(j)}x_{j,\mathrm{raw}}.
\end{equation}

B 和 S 直接使用限制到 $[0,1]$ 的分项结果质量；视觉引导分域还加入更保守的条件：

\begin{align}
q_E &= \min\left(q_{\mathrm{result}},q_{\mathrm{trial}},
\operatorname{clip}\left(\frac{\mathrm{SNR}_{\min}}{3.5}\right),
\operatorname{clip}\left(1-\frac{e_{\mathrm{repeat}}}{1.0}\right)\right),\\
q_F &= \min\left(q_{\mathrm{result}},
\operatorname{clip}\left(\frac{s_{\mathrm{smile}}}{0.22}\right)\right),\\
q_A &= \min\left(q_{\mathrm{result}},q_{\mathrm{completion}}\right).
\end{align}

其中 \texttt{clip} 表示把数值限制在 $[0,1]$。这样会降低边缘质量条件下特征的幅度，但仍保留对应的质量值和缺失信息，供后续模型显式使用。

\subsection{可解释分域严重度}

除原始和门控向量外，融合层还为每个可用分域计算相对工程阈值的归一化严重度。B、E、F、A 使用各自规则中相同的研究参考值：

\begin{align}
r_B &= \max\left(\frac{z_{\mathrm{trunk}}}{3.5},\frac{z_{\mathrm{sway}}}{3.5}\right),\\
r_E &= \max\left(\frac{a_{\mathrm{direction}}}{0.45},
\frac{e_{\mathrm{conjugacy}}}{0.35},\frac{d_{\mathrm{rest}}}{12}\right),\\
r_F &= \max\left(\frac{|\Delta d_{\mathrm{mouth}}|}{0.075},
\frac{|\Delta s_{\mathrm{smile}}|}{0.24}\right),\\
r_A &= \max\left(\frac{|\Delta h|}{0.30},
\frac{|\Delta_{\mathrm{drift}}|}{0.22}\right).
\end{align}

S 使用字符错误率相对 0.35、停顿比例相对 0.55，以及偏离每秒 1--8 个字符参考范围的语速变化。最终发布的分域值为：

\begin{equation}
d_c=q_c\min(5,r_c),
\end{equation}

其中 5 是可配置的截断值，用于避免极端比值支配探索性汇总。这些数值只表示相对工程阈值的幅度，不是经过校准的概率或临床严重程度。

\subsection{聚合描述量}

融合载荷输出非零质量分域数量、最大分域严重度，以及两个最高可用分域严重度的平均值。同时输出 A/F 侧别一致性：当 A 和 F 严重度均非零且主要有符号特征均非零时，同号输出 $+1$，异号输出 $-1$，不可用或无方向证据输出 0。B 和 E 不参与该侧别描述，因为其画面方向不能合理推断病灶侧别。

这些聚合量仅用于描述。任何聚合值阈值都没有接入用户可见的筛查决策。

\subsection{测试与未来模型使用}

专项测试覆盖特征顺序和向量长度、质量乘法、各分项质量规则、缺失掩码、归一化严重度、A/F 一致性，以及把仅人工报告结果视为传感器缺失。未来如果使用该 28 维向量训练预测模型，必须按受试者划分数据，固定并版本化预处理及特征顺序，研究缺失机制，完成校准和独立外部验证，并在每条训练样本及推理输出中保存融合版本。

\section{模块集成、报告与隐私}

\subsection{会话状态与独立检测}

线程安全的会话控制器为每个模块保存独立检测器，并区分待机、准备、采集、重试和报告状态。各项目可按任意顺序执行。重复检测会创建新的编号报告，而不是覆盖前一次结果。采集过程中切换摄像头会重新开始当前项目，防止不同视角的坐标进入同一次判定。每次会话状态快照都会在普通分项报告旁构建并暴露当前融合载荷。

\subsection{人工与自动证据融合}

B 和 E 中，已完成的人工观察如果报告异常，在报告中优先于自动检测。F 和 A 来自相应视觉模块，S 来自麦克风和离线识别器。报告保存状态、原因、来源、质量、量化指标，以及可选的侧别和详细信息。界面应提供“左臂持续较低”“眼动响应不可重复”或“录音音量过低”等可理解说明，而不是只显示无解释的分数。

\subsection{隐私与故障行为}

推理运行在本地服务主机，浏览器不需要向第三方云端发送数据。个人平衡和被动语音基线保存在本机，日常被动语音 WAV 在分析后删除。研究日志和阳性事件记录均可配置。模型、识别器或麦克风不可用时，相应项目返回能力不足，不得伪造阴性结果。

\section{测试与验证方案}

\subsection{软件验证}

项目包含 BE-FAST 几何与判定、会话历史、监测、Web API、预览、训练工具和故障路径的自动测试。软件验证应覆盖：

\begin{itemize}
  \item 关键点几何和坐标符号；
  \item 阈值边界条件；
  \item 缺失、低置信度和遮挡关键点；
  \item 言语质量门控、转写归一化、CER、停顿和语速边界；
  \item 被动语音基线持久化、多域投票、暂停/恢复和 WAV 删除；
  \item 阶段转换、重试、跳过和摄像头切换；
  \item 指标及原因码的序列化；
  \item 融合向量顺序、质量门控、缺失掩码、严重度和侧别描述量；
  \item 可选模型或硬件缺失与失败。
\end{itemize}

2026 年 8 月实现快照中，仓库 128 项自动测试全部通过，包括共享语音表征提取、模型加载、说话人无泄漏 manifest 生成和影子失败隔离。通过上述测试只能说明软件符合编码规格，不能证明临床准确性。

\subsection{工程性能评估}

应在目标树莓派配置上测量有效推理帧率、端到端延迟、CPU、内存、温度、丢帧率和摄像头故障恢复。应在不同摄像头、距离、分辨率、光照、眼镜、衣物和体型条件下测试重复性。每个模块都应报告数据不足率和重测成功率。

\subsection{临床研究验证}

在把自动阳性用作研究终点前，前瞻性方案应纳入急性脑卒中、卒中类似疾病、相关既往功能障碍和健康对照。训练集、开发集与测试集应按受试者划分，而不是按视频帧或音频窗口划分。E 应与同步临床或研究级眼动仪比较；B 应与适当的平衡仪器和临床 lateropulsion 评估比较；F 和 A 应与盲法临床评分及标准化视频标注比较；S 应在真实声学条件下与盲法言语语言评估比较。开发阈值应在独立验证前冻结，并报告灵敏度、特异度、ROC、PR、置信区间、校准度、评分者一致性、重复性和数据不足比例。

筛查用途通常偏向灵敏度，但这并不意味着可以在缺乏代表性数据时任意选择阈值。正式研究还需要预先定义不确定灰区，并保持“突发症状立即升级处理”的规则。

\section{风险、局限与安全控制}

主要风险包括错误安心、误报警、基线不具代表性、动作理解错误以及不同设备和人群间性能差异。当前控制措施包括症状优先、标准化动作、画面内引导、重复测量、显式质量门控、保守的数据不足状态、本地处理和明确的急症提示。这些设计可以降低可预见的工程故障，但不能证明医疗安全或监管合规。

最重要的解释边界是：\textbf{阴性只表示在质量合格的任务中没有检测到目标可见表型，并不表示没有脑卒中。}同样，阳性表示警示症状或工程定义表型，并不确认病因。

\section{总结与后续工作}

PiBE-FAST 将 B、E、F、A、S 实现为透明的时序测量，而不是单帧或单分数诊断。B 使用个人稳健基线比较躯干方向和横向运动；E 结合重复配对终点、头动和噪声门控；F 将微笑运动与中性面部基线比较；A 在评估高度和下落前要求完整完成双臂动作；S 则把持续个人声学变化与本地转写固定句确认分开处理。主动 S 报告现在还以影子模式携带经过说话人无泄漏内部评估的 MDSC 构音障碍表征。融合层在不改变安全决策的前提下，把已完成的连续测量保存在质量门控的 28 维研究表示中。阶段化调度为进一步开展树莓派评估提供了基础，结构化结果状态与显式缺失掩码则保留了测量不确定性。

后续工作应建立设备一致的采集协议，采集经过伦理治理和临床标注的数据，量化不同人群和硬件的性能，重新标定全部研究阈值，并进行前瞻性完整流程验证。在此之前，系统仍是支持研究和演示的原型，不能替代紧急专业评估。

\bibliographystyle{unsrt}
\bibliography{references}

\appendix
\section{BEFA 默认参数}

\begin{longtable}{p{0.22\linewidth}p{0.31\linewidth}p{0.37\linewidth}}
\caption{报告快照对应的实现默认值。}\\
\toprule
模块 & 参数 & 默认值及作用 \\
\midrule
\endfirsthead
\toprule
模块 & 参数 & 默认值及作用 \\
\midrule
\endhead
公共 & MoveNet 最低分数 & 0.35，可见性门控 \\
B & 准备期/采集期 & 1.5 秒/30 秒 \\
B & 有效数据 & 至少 30 个样本且有效率不低于 0.75 \\
B & 个人基线 & 5 个合格窗口 \\
B & 稳健变化分数 & 3.5，仅用于变化监测 \\
E & 单目标/稳定时间 & 2.0 秒/0.5 秒 \\
E & 重复次数 & 每个侧方向 3 次 \\
E & 有效数据 & 每试次至少 5 个样本且有效率不低于 0.60 \\
E & 眼宽/眼位 MAD & 至少 24 像素/不高于 0.08 \\
E & 头姿覆盖/变化 & 至少 0.80/不超过 $8^{\circ}$ \\
E & 响应 SNR & 3.5 \\
E & 静息/方向/共轭 & $12^{\circ}$/0.45/0.35，研究阈值 \\
F & 中性/微笑 & 2 秒/3 秒 \\
F & 有效数据 & 每阶段至少 5 个样本且总体有效率不低于 0.50 \\
F & 微笑强度 & 至少 0.22 \\
F & 几何差/blendshape 差 & 0.075/0.24 \\
A & 预览/自动准备/倒计时 & 2.0 秒/1.2 秒/3.0 秒 \\
A & 抬起/保持/放下 & 8 秒/5 秒/5 秒 \\
A & 抬臂姿势 & 手腕容差 0.42、肘角 $\geq125^{\circ}$、伸展 $\geq0.45$ \\
A & 有效保持样本 & 至少 20 个；名义有效率 0.55 \\
A & 高度差/下落差 & 0.30/0.22 个肩宽归一化单位 \\
S 被动 & 窗口/间隔 & 8 秒/1 秒 \\
S 被动 & 音频质量 & 时长 $\geq2$ 秒、发声 $\geq1.5$ 秒、RMS $\geq-42$ dBFS、削波 $\leq0.02$ \\
S 被动 & 基线/变化规则 & 6 个窗口；至少 2 个域中 $z\geq3.5$ \\
S 被动 & 持续投票 & 最近 5 个有效窗口至少 3 票异常 \\
S 主动 & 采集格式 & 约 7 秒、16 kHz、16-bit 单声道 PCM \\
S 主动 & 音频质量 & 时长 $\geq2$ 秒、发声 $\geq1$ 秒、RMS $\geq-42$ dBFS、削波 $\leq0.02$ \\
S 主动 & CER/停顿 & 0.35/0.55 \\
S 主动 & 语速范围 & 每秒 1--8 个识别字符 \\
融合 & 连续特征/模型维数 & 18/28 \\
融合 & 分项顺序 & B、E、F、A、S \\
融合 & 严重度截断值 & 5.0，仅作研究描述 \\
融合 & S 参考值 & CER 0.35、停顿 0.55、语速每秒 1--8 字符 \\
\bottomrule
\end{longtable}

\section{主要源文件}

\begin{longtable}{p{0.42\linewidth}p{0.48\linewidth}}
\toprule
文件 & 作用 \\
\midrule
\path{app/befast/config.py} & 时间与阈值配置 \\
\path{app/befast/balance.py} & B 特征、个人基线与判定 \\
\path{app/befast/eyes.py} & E 试次、质量门控与判定 \\
\path{app/befast/face.py} & F 中性到微笑比较 \\
\path{app/befast/arms.py} & A 动作阶段与时序特征 \\
\path{app/befast/speech.py} & S 统一报告映射 \\
\path{app/speech_audio.py} & 主动录音、ASR 与评估 \\
\path{app/passive_speech.py} & 自然语音个人基线监测 \\
\path{app/befast/fusion.py} & 质量门控固定长度研究向量 \\
\path{app/befast/session.py} & 线程安全流程编排 \\
\path{app/befast/report.py} & 人工与自动结果组装 \\
\path{tests/test_befast.py} & 模块与会话验证 \\
\path{tests/test_fusion.py} & 融合向量与缺失语义验证 \\
\path{tests/test_speech_audio.py} & 主动 S 音频与 ASR 验证 \\
\path{tests/test_passive_speech.py} & 被动 S 特征与监测器验证 \\
\bottomrule
\end{longtable}

\end{document}
